\documentclass[UTF8,11pt]{ctexart}
\usepackage[a4paper,margin=24mm]{geometry}
\usepackage{amsmath,amssymb,amsthm,mathtools,mathrsfs}
\usepackage[all]{xy}
\usepackage{xurl,hyperref,enumitem}
\usepackage{tikz}
\usetikzlibrary{arrows.meta,cd,decorations.markings,calc,patterns}
\hypersetup{hidelinks,breaklinks=true}
\setlength{\emergencystretch}{3em}
\allowdisplaybreaks
\newtheorem*{theorem}{Theorem}
\newtheorem*{thm}{Theorem}
\newtheorem*{lemma}{Lemma}
\newtheorem*{proposition}{Proposition}
\newtheorem*{prop}{Proposition}
\newtheorem*{corollary}{Corollary}
\newtheorem*{cor}{Corollary}
\newtheorem*{claim}{Claim}
\theoremstyle{definition}
\newtheorem*{definition}{Definition}
\newtheorem*{defn}{Definition}
\newtheorem*{remark}{Remark}
\newtheorem*{example}{Example}
\newcommand{\R}{\mathbb R}
\newcommand{\C}{\mathbb C}
\newcommand{\Z}{\mathbb Z}
\newcommand{\Q}{\mathbb Q}
\newcommand{\N}{\mathbb N}
\newcommand{\F}{\mathbb F}
\newcommand{\D}{\mathbb D}
\newcommand{\bB}{\mathbb B}
\newcommand{\bC}{\mathbb C}
\newcommand{\bR}{\mathbb R}
\newcommand{\bZ}{\mathbb Z}
\newcommand{\bN}{\mathbb N}
\newcommand{\bQ}{\mathbb Q}
\newcommand{\bD}{\mathbb D}
\newcommand{\bF}{\mathbb F}
\newcommand{\CP}{\mathbb{CP}}
\newcommand{\RP}{\mathbb{RP}}
\newcommand{\sO}{\mathscr O}
\newcommand{\sI}{\mathscr I}
\newcommand{\sF}{\mathscr F}
\newcommand{\sG}{\mathscr G}
\newcommand{\sH}{\mathscr H}
\newcommand{\sR}{\mathscr R}
\newcommand{\sM}{\mathscr M}
\newcommand{\sV}{\mathscr V}
\newcommand{\sU}{\mathscr U}
\DeclareMathOperator{\Hom}{Hom}
\DeclareMathOperator{\Ext}{Ext}
\DeclareMathOperator{\Tor}{Tor}
\DeclareMathOperator{\Spec}{Spec}
\DeclareMathOperator{\Proj}{Proj}
\DeclareMathOperator{\supp}{supp}
\DeclareMathOperator{\im}{im}
\DeclareMathOperator{\id}{id}
\DeclareMathOperator{\Tr}{Tr}
\DeclareMathOperator{\tr}{tr}
\DeclareMathOperator{\rank}{rank}
\DeclareMathOperator{\ord}{ord}
\DeclareMathOperator{\dist}{dist}
\DeclareMathOperator{\End}{End}
\DeclareMathOperator{\Aut}{Aut}
\DeclareMathOperator{\Gal}{Gal}
\DeclareMathOperator{\vol}{vol}
\DeclareMathOperator{\Cl}{Cl}
\DeclareMathOperator{\coker}{coker}
\DeclareMathOperator{\codim}{codim}
\DeclareMathOperator{\divergence}{div}
\DeclareMathOperator{\Ric}{Ric}
\DeclareMathOperator{\grad}{grad}
\DeclareMathOperator{\Hess}{Hess}
\newcommand{\abs}[1]{\left\lvert#1\right\rvert}
\newcommand{\sabs}[1]{\lvert#1\rvert}
\newcommand{\babs}[1]{\bigl\lvert#1\bigr\rvert}
\newcommand{\Babs}[1]{\Bigl\lvert#1\Bigr\rvert}
\newcommand{\bbabs}[1]{\biggl\lvert#1\biggr\rvert}
\newcommand{\BBabs}[1]{\Biggl\lvert#1\Biggr\rvert}
\newcommand{\norm}[1]{\left\lVert#1\right\rVert}
\newcommand{\snorm}[1]{\lVert#1\rVert}
\newcommand{\bnorm}[1]{\bigl\lVert#1\bigr\rVert}
\newcommand{\Bnorm}[1]{\Bigl\lVert#1\Bigr\rVert}
\newcommand{\bbnorm}[1]{\biggl\lVert#1\biggr\rVert}
\newcommand{\BBnorm}[1]{\Biggl\lVert#1\Biggr\rVert}
\newcommand{\widebar}[1]{\overline{#1}}
\newcommand{\nicefrac}[2]{\frac{#1}{#2}}
\newcommand{\linnprod}[2]{\langle#1,#2\rangle}
\newcommand{\myindex}[1]{#1}
\newcommand{\glsadd}[1]{}
\newcommand{\avoidbreak}{}
\newcommand{\sourceref}[1]{\textnormal{[原文：\nolinkurl{#1}]}}
\newcommand{\thmref}[1]{\sourceref{#1}}
\newcommand{\lemmaref}[1]{\sourceref{#1}}
\newcommand{\propref}[1]{\sourceref{#1}}
\newcommand{\corref}[1]{\sourceref{#1}}
\newcommand{\defnref}[1]{\sourceref{#1}}
\newcommand{\exampleref}[1]{\sourceref{#1}}
\newcommand{\exerciseref}[1]{\sourceref{#1}}
\newcommand{\figureref}[1]{\sourceref{#1}}
\newcommand{\sectionref}[1]{\sourceref{#1}}
\newcommand{\appendixref}[1]{\sourceref{#1}}
\newcommand{\stacksref}[2]{\href{https://stacks.math.columbia.edu/tag/#1}{#2 [#1]}}

\begin{document}
\section*{031\quad 全纯函数的Weierstrass预备分解}
\subsection*{来源与计数目标}
Ji\v{r}\'i Lebl, \emph{Tasty Bits of Several Complex Variables}。从作者公开的原生 TeX 正文摘录；获取日期：2026-09-28。使用实际访问的在线正文，不把工作分支冒称冻结的印刷版。
\par\url{https://raw.githubusercontent.com/jirilebl/scv/master/scv.tex}
\par\url{https://www.jirka.org/scv/}
\par 原文定位：Chapter 6，\texttt{sec:wpt}，Weierstrass preparation theorem 及全部证明；前置 Weierstrass polynomial 定义。
\par\textbf{本文件只计一个问题：}全纯函数分解为非消失全纯单位与唯一 Weierstrass 多项式；不把相邻除法定理另计入本题。
\subsection*{作者命题与人类证明}
\begin{definition}[Weierstrass polynomial]
Let $U\subset\C^{n-1}$ be open, $0\in U$, and let $z'\in\C^{n-1}$ denote the coordinates. Suppose $P\in\sO(U)[z_n]$ is a monic polynomial of degree $k\geq0$,
\[
P(z',z_n)=z_n^k+\sum_{\ell=0}^{k-1}c_\ell(z')z_n^\ell,
\]
where $c_\ell$ are holomorphic functions defined on $U$ such that $c_\ell(0)=0$ for all $\ell$. Then $P$ is called a Weierstrass polynomial of degree $k$.
\end{definition}
\begin{theorem}[Weierstrass preparation theorem]
Suppose $f\in\sO(U)$ for an open $U\subset\C^{n-1}\times\C$, where $0\in U$, and $f(0)=0$. Suppose $z_n\mapsto f(0,z_n)$ is not identically zero near the origin and its order of vanishing at the origin is $k\geq1$. Then there exists an open polydisc $V=V'\times D\subset\C^{n-1}\times\C$ with $0\in V\subset U$, a unique $u\in\sO(V)$, $u(z)\neq0$ for all $z\in V$, and a unique Weierstrass polynomial $P$ of degree $k$ with coefficients holomorphic in $V'$ such that
\[
f(z',z_n)=u(z',z_n)P(z',z_n),
\]
and such that all $k$ zeros (counting multiplicity) of $z_n\mapsto P(z',z_n)$ lie in $D$ for all $z'\in V'$.
\end{theorem}
\begin{proof}
There exists a small disc $D\subset\C$ centered at zero such that $\{0\}\times\widebar D\subset U$ and such that $f(0,z_n)\neq0$ for $z_n\in\widebar D\setminus\{0\}$. By continuity of $f$, there is a small polydisc $V=V'\times D$ such that $\widebar V\subset U$ and $f$ is not zero on $V'\times\partial D$. See \figureref{fig:wpreppoly} for the setup. We will consider the zeros of $z_n\mapsto f(z',z_n)$ for $z'\in V'$. See \figureref{fig:wprepdiv}.

By the one-variable argument principle (\thmref{thm:onevarargprinc}), the number of zeros (with multiplicity) of $z_n\mapsto f(z',z_n)$ in $D$ is
\[
\frac1{2\pi i}\int_{\partial D}\frac{\frac{\partial f}{\partial z_n}(z',\zeta)}{f(z',\zeta)}\,d\zeta.
\]
As $f(z',\zeta)$ does not vanish when $z'\in V'$ and $\zeta\in\partial D$, the expression above is a continuous integer-valued function of $z'\in V'$. The expression is equal to $k$ when $z'=0$, and so it is equal to $k$ for all $z'\in V'$.
Write the zeros of $z_n\mapsto f(z',z_n)$ as $\alpha_1(z'),\ldots,\alpha_k(z')$, including multiplicity. The zeros are not ordered in any particular way---pick \emph{some} ordering for every $z'$. Write
\[
P(z',z_n)=\prod_{\ell=1}^k\bigl(z_n-\alpha_\ell(z')\bigr)=z_n^k+c_{k-1}(z')z_n^{k-1}+\cdots+c_0(z').
\]
For a fixed $z'$, $P$ (and thus the coefficients $c_0,\ldots,c_{k-1}$) is uniquely defined as its definition is independent of the ordering of the zeros. That $c_j(0)=0$ for all $j$ follows as $\alpha_\ell(0)=0$ for all $\ell$. As the above is the unique way to define a monic polynomial with these zeros (\exerciseref{exercise:monicpolyunique}), the uniqueness part of the theorem follows.

We need to show that the coefficients are holomorphic functions on $V'$, and that $u$ is a holomorphic function on $V$. No matter how you ordered the zeros for each $z'$, the functions $\alpha_\ell$ may not be continuous in general (see \exampleref{sqrt:example}). However, we will prove that the functions $c_\ell$ are holomorphic. The functions $c_\ell$ are (up to sign) the elementary symmetric functions of $\alpha_1,\ldots,\alpha_k$ (see below). A standard theorem in algebra (Newton's identities, see \exerciseref{exercise:powersums}) says that the elementary symmetric functions are polynomials in the so-called power sum functions in the $\alpha_\ell$s:
\[
s_m(z')=\sum_{\ell=1}^k\alpha_\ell(z')^m,\qquad m=1,\ldots,k.
\]
So if the power sums $s_m$ are holomorphic on $V'$, then $c_\ell$ are holomorphic on $V'$. A refinement of the argument principle (see \thmref{thm:onevarargprinc}) says: If $h$ and $g$ are holomorphic functions on a disc $D$, continuous on $\widebar D$, such that $g$ has no zeros on $\partial D$, and $\alpha_1,\ldots,\alpha_k$ are the zeros of $g$ in $D$, then
\[
\frac1{2\pi i}\int_{\partial D}h(\zeta)\frac{g'(\zeta)}{g(\zeta)}\,d\zeta=\sum_{\ell=1}^kh(\alpha_\ell).
\]
With $h(\zeta)=\zeta^m$ and $g(\zeta)=f(z',\zeta)$, the theorem says
\[
s_m(z')=\sum_{\ell=1}^k\alpha_\ell(z')^m=\frac1{2\pi i}\int_{\partial D}\zeta^m\frac{\frac{\partial f}{\partial z_n}(z',\zeta)}{f(z',\zeta)}\,d\zeta.
\]
The function $s_m$ is clearly continuous, and if we differentiate under the integral with $\frac\partial{\partial\bar z_1},\ldots,\frac\partial{\partial\bar z_{n-1}}$ we find that $s_m$ is holomorphic. Thus $c_0,\ldots,c_{k-1}$ are holomorphic, and so $P$ is a Weierstrass polynomial.

Finally, we wish to show that $P$ divides $f$ as claimed, that is, that $u$ is holomorphic. For each fixed $z'$, one-variable theory says that $z_n\mapsto\frac{f(z',z_n)}{P(z',z_n)}$ has only removable singularities, and in fact, it has no zeros as we defined $P$ to exactly cancel them all out. The Cauchy formula on $\nicefrac fP$ then says that the function
\[
u(z',z_n)=\frac1{2\pi i}\int_{\partial D}\frac{f(z',\zeta)}{P(z',\zeta)(\zeta-z_n)}\,d\zeta
\]
is equal to $\frac{f(z',z_n)}{P(z',z_n)}$. The function $u$ is clearly continuous and holomorphic in $z_n$ for each fixed $z'$. Differentiating under the integral shows it is also holomorphic in $z'$.
\end{proof}

\subsection*{整理说明（非作者正文）}
$\sO(U)$ 表示 $U$ 上的全纯函数环。允许的标准先修为一变量辐角原理及其带权留数形式、Cauchy 公式、可去奇点、Newton 恒等式，以及参数积分求导。核心工作是处理无法连续编号的移动零点，以幂和积分构造全纯系数，再构造非消失商函数；这些均保留，不作为先修略去。

原文两幅图仅图示已在文字中完整定义的多圆柱邻域和移动零点，不承担额外推理；本摘录保留原图标签供回查，不收入图像文件。外部定理、习题和图的引用用原 TeX 标签显示，均指原书而非本文件编号。源书将 Newton 恒等式列作练习，本题明确把这项标准代数事实作为先修，不冒称已收其证明。原证明的 ``see below'' 指原书随后对初等对称函数的定义。保留原语言和段落论证顺序，仅规范公式空白、宏与断行。
\subsection*{可修改的简短标注（非作者正文）}
$c$：全纯函数局部零点与因子分解

$a$：辐角原理；留数积分；移动零点；初等对称函数；Newton 恒等式；参数积分；Cauchy 公式

\texttt{cross\_theory\_status = yes}。把多值零点数据经对称幂和变成单值全纯系数，以代数多项式编码解析零集，再用 Cauchy 积分恢复单位因子。位置：原证明的 $s_m$ 积分和最后的 $u$ 积分。

全部标注均为非穷尽、可修改初稿；未标注不表示未使用。
\subsection*{许可与修改记录}
原数学正文作者：Ji\v{r}\'i Lebl。作者网站及源书声明允许 CC BY-SA 4.0 或 CC BY-NC-SA 4.0 双许可；本条选择 CC BY-SA 4.0，整理部分亦采用同一许可。\url{https://creativecommons.org/licenses/by-sa/4.0/}。2026-09-28：助手截取题证、调整独立排版，加入来源、范围与初稿标注；不暗示作者认可本次整理。保留的是所用 TeX 正文摘录，未将其称为整书原件。
\end{document}
